\section{Introducción}

El análisis computacional de salud mental en texto generado por usuarios se ha consolidado como una línea de investigación estratégica para la detección temprana de señales de riesgo, especialmente en entornos donde la búsqueda de ayuda profesional se retrasa por estigma, costo o falta de acceso. En ese contexto, las plataformas de redes sociales ofrecen narrativas espontáneas, longitudinales y, en muchos casos, suficientemente detalladas para estudiar indicadores de depresión, desesperanza, aislamiento, autolesión e ideación suicida \citep{montejo2024survey, abdulsalam2024suicidal}.

Dentro de ese ecosistema, Reddit ocupa una posición particularmente relevante por su estructura basada en comunidades temáticas, la extensión relativamente mayor de sus publicaciones y la posibilidad de anonimato, factores que favorecen expresiones más explícitas sobre malestar psicológico. Trabajos tempranos y recientes muestran que el valor de Reddit no se limita a la clasificación binaria de riesgo, sino que también permite modelar transiciones hacia ideación suicida, niveles de severidad, categorías diagnósticas clínicamente significativas y explicaciones textuales de apoyo a la decisión \citep{dechoudhury2016shifts, gaur2019severity, bao2025redsm5}.

Este documento presenta una revisión de estado del arte, no un artículo experimental. Su objetivo es identificar y analizar literatura reciente y representativa sobre detección computacional de señales de salud mental en redes sociales, con énfasis en depresión, ideación suicida, marcos DSM-5, aprendizaje automático, aprendizaje profundo, modelos basados en transformers y modelos de lenguaje de gran escala. La revisión se organiza en función de la dimensión de desempeño seleccionada para el reto: \textit{efectividad predictiva e interpretabilidad clínica}. Esta elección permite evaluar no solo qué tan bien predicen los modelos, sino también si sus decisiones pueden vincularse con síntomas, criterios o evidencias comprensibles para profesionales de salud mental.

La hipótesis que articula la revisión es que el campo ha avanzado con rapidez en capacidad predictiva, pero no de manera uniforme en fundamentación clínica, transparencia y utilidad para intervención temprana. Por ello, el documento compara enfoques centrados en características lingüísticas, modelos profundos, recursos clínicamente anotados y sistemas basados en transformers y LLMs, destacando sus aportaciones, limitaciones y el tipo de evidencia que ofrecen para fases posteriores del proyecto.
